\begin{table}[H]
\centering
\caption{Sensitivity to Delivery-Fidelity Calibration}
\label{tab:struct_tau_sensitivity}
\begin{threeparttable}
\footnotesize
\begin{tabular}[t]{lcccccc}
\toprule
Spec. & $\tau_K$ & $\tau_L$ & $\tau_N$ & Max ITT/SE & K high-$\tau$ & High-input \\
\midrule
Lower process map & 0.40 & 0.31 & 0.22 & 0.297 & $+0.179$ & $+0.206$ \\
Preferred & 0.50 & 0.38 & 0.31 & 0.281 & $+0.155$ & $+0.187$ \\
Upper process map & 0.60 & 0.46 & 0.40 & 0.276 & $+0.126$ & $+0.165$ \\
\bottomrule
\end{tabular}
\begin{tablenotes}[para,flushleft]
\footnotesize
\item \textit{Notes:} Each row re-estimates the stage-4 production block after replacing the stage-3 delivery-fidelity primitives with a transparent lower, preferred, or upper process-to-$\tau$ mapping. The lower and upper mappings preserve the same process moments and market ordering logic but vary the intercept and slope used to translate completion evidence for the assignment channel into $\tau$. The high-input counterfactual still sets $\rho=\max_m\hat{\rho}_m$, $\omega=0.98$, and $\tau=0.95$; this table asks whether the observed-cell calibration of $\tau$ materially changes the extrapolation.
\end{tablenotes}
\end{threeparttable}
\end{table}
